{\actuality} Текстовые документы остаются основным средством хранения
и~передачи информации в~деловой, учебной и~научной деятельности. С~развитием
цифровых каналов коммуникации объём текстов, с~которым сталкивается отдельный
пользователь или организация, увеличивается многократно: отчёты, заметки,
статьи, переписка, техническая документация, уведомления и~согласования
накапливаются в~личных и~корпоративных хранилищах быстрее, чем поддаются
ручному упорядочению. Это делает актуальной задачу автоматической
классификации текстовых документов "--- их~распределения по~тематическим,
проектным или функциональным категориям, отражающим практику конкретного
пользователя.

Среди текстовых потоков, требующих регулярной классификации, особое место
занимает электронная почта. Она остаётся одним из~основных каналов деловой,
учебной и~научной коммуникации, и~для профессионального пользователя почтовый
ящик постепенно превращается в~хранилище задач, уведомлений, документов,
согласований и~технических сообщений, объём которого может достигать
нескольких сотен входящих писем в~день \autocite{Singh2023EmFORE}. Письма
одновременно сочетают признаки служебной переписки, отчётности, технических
уведомлений и~приложенных файлов, что делает их~характерным и~достаточно
сложным частным случаем общей задачи классификации текстов. Ручная сортировка
такой корреспонденции по~папкам требует регулярных временны́х затрат, а~набор
простых правил по~отправителю, теме или отдельным словам быстро устаревает:
меняются проекты, контрагенты, форматы уведомлений и~язык переписки.

Готовые почтовые клиенты предлагают встроенные механизмы организации входящих
сообщений "--- например, Focused~Inbox в~Outlook и~Priority~Inbox в~Gmail.
Однако подобные эвристики ориентированы на~общие категории «важное~/ неважное»
и~не~учитывают индивидуальную структуру почтового архива конкретного
пользователя. В~настоящей работе задача понимается шире: система должна
автоматически относить сообщение к~персонализированному классу "--- проекту,
направлению работы, типу уведомления или иной категории, принятой самим
пользователем.

Отдельное значение имеет локальный характер обработки данных. Тексты писем
могут содержать персональные данные, коммерческую информацию, вложения,
внутренние обсуждения и фрагменты документов. Поэтому передача содержимого
писем во~внешние сервисы не~всегда допустима из-за регламентов обработки
данных и~внутренних правил организации. В~настоящей работе рассматривается
локальная система, в~которой синхронизация, нормализация, построение векторных
представлений текста, кластеризация и~классификация выполняются на~оборудовании
пользователя или в~контролируемом им окружении.

Почтовые архивы часто смешивают русско- и~англоязычные сообщения, технические
уведомления, короткие ответы, длинные цепочки переписки, подписи, цитаты
предыдущих писем и~HTML-разметку. Это требует устойчивой предобработки текста
и~multilingual-моделей, способных представлять сообщения разных языков
в~едином признаковом пространстве.

{\progress} Задача автоматической классификации текстовых документов
исследуется в~рамках обработки естественного языка с~момента появления
методов статистического анализа текстов и~охватывает широкий круг
прикладных сценариев: рубрицирование новостей, тематическое моделирование
научных публикаций, фильтрацию спама, упорядочение служебной переписки.
Применительно к~электронной почте задача имеет собственную историю
исследований. В~работе V.~L.~Keiser она рассматривается на~примере
предсказания проектов в~системе TaskTracer; сравниваются наивные байесовские
классификаторы, методы Passive-Aggressive и~Confidence-Weighted Linear
Classification для онлайн-классификации писем
\autocite{Keiser2009TaskTracer}. Современный подход EmFORE предлагает
символьно-правиловое онлайн-обучение правил распределения писем по~папкам
и~отдельно рассматривает изменение пространства признаков и~дрейф
пользовательских категорий \autocite{Singh2023EmFORE}. Эти работы показывают,
что задача не~сводится к~разовой классификации фиксированного набора текстов:
система должна учитывать персональные классы, появление новых типов сообщений,
низкоуверенные решения и~обратную связь пользователя.

В~настоящей работе рассматривается следующий вариант решения: применить
локальные multilingual-эмбеддинги, выполнить кластеризацию исторической почты
для начального формирования классов и~обучить supervised-классификатор
на~полученной разметке. Такой подход позволяет не~задавать категории вручную
до~начала эксплуатации, оставить текст писем внутри локального контура
и~выделить сообщения вне известных классов в~отдельный поток для проверки.

{\aim} настоящей работы является разработка интеллектуальной системы
классификации текстовых документов на~примере электронной почты, способной
работать локально на~оборудовании пользователя, поддерживать
персонализированные классы, формировать их по~исторической выборке
и~оценивать качество принимаемых решений.

Для достижения поставленной цели необходимо решить следующие {\tasks}:
\begin{enumerate}[beginpenalty=10000]
  \item проанализировать существующие подходы к~классификации текстовых
        документов и~автоматической организации электронной почты;
  \item сформулировать функциональные и~нефункциональные требования
        к~локальной системе классификации писем;
  \item спроектировать сервисно-ориентированную архитектуру системы и~модель
        данных для хранения сырых писем, нормализованных текстов, эмбеддингов,
        меток, предсказаний и~версий моделей;
  \item разработать алгоритм автоматического формирования первичных классов
        на~основе локальных multilingual-эмбеддингов и~кластеризации
        исторической выборки;
  \item реализовать supervised-классификатор текстовых сообщений на~базе
        локального представления текста и~обучаемой классификационной
        надстройки;
  \item реализовать механизм принятия решений с~явной категорией
        \texttt{other}, порогом уверенности, сохранением версии модели
        и~объяснимых атрибутов предсказания;
  \item программно реализовать MVP-систему и~провести экспериментальное
        исследование качества классификации, кластеризации и~поведения
        на~сообщениях вне известных классов;
  \item сформулировать критерии активации модели и~стратегию обратной связи
        через повторную аннотацию сообщений, попавших в~категорию
        \texttt{other}.
\end{enumerate}

\textbf{Объект исследования.} Процесс автоматической многоклассовой
классификации текстовых документов с~учётом персонализированных
пользовательских классов.

\textbf{Предмет исследования.} Методы и~алгоритмы построения локальной
системы классификации электронной почты, включая нормализацию текста,
кластеризацию для начального формирования классов, обучение классификатора
и~обработку сообщений, не~относящихся к~известным категориям.

{\methods} В~работе применяются методы системного анализа предметной области,
методы обработки естественного языка, методы машинного обучения для
классификации текстов, алгоритмы построения векторных представлений
предложений, методы кластеризации на~основе плотности, линейные и~нейросетевые
классификаторы. Для оценки результатов используются accuracy, macro-F1,
precision и~recall по~классам, confusion matrix, coverage, доля сообщений,
отнесённых к~\texttt{other}, и~операционные показатели обработки почтового
потока. Воспроизводимость экспериментов обеспечивается фиксированной
конфигурацией, сохранением артефактов моделей и~версионированием наборов
данных.

{\novelty}
\begin{enumerate}[beginpenalty=10000]
  \item Предложен локальный pipeline классификации электронной почты,
        объединяющий нормализацию MIME-сообщений, построение
        multilingual-эмбеддингов, кластеризацию исторической выборки
        и~supervised-классификацию новых сообщений без~передачи текста
        во~внешние сервисы.
  \item Разработан способ начального формирования пользовательских классов
        через кластеризацию исторической почты с~выделением outliers
        и~последующей ручной проверкой кандидатов.
  \item Предложена схема принятия решений с~явной категорией \texttt{other},
        при которой низкоуверенные и~потенциально out-of-distribution
        сообщения сохраняются как материал для обратной связи и~последующего
        переобучения.
  \item Сформулированы quality gates для активации модели: совместная
        проверка качества кластеризации, качества классификатора
        и~операционных характеристик pipeline до~включения автоматических
        действий с~почтовым ящиком.
\end{enumerate}

{\influence} работы состоит в~создании MVP-прототипа локальной системы
классификации электронной почты. Прототип может использоваться как основа
для развёртывания в~пользовательском окружении, где требуется сохранить
приватность почтовых данных и~при этом получить персонализированную
автоматическую классификацию. Практический результат включает программные
компоненты read-only IMAP-синхронизации, нормализации писем, кластеризации,
обучения модели, инференса и~сохранения предсказаний в~базе данных
без~изменения состояния почтового ящика.

{\defpositions}
\begin{enumerate}[beginpenalty=10000]
  \item Локальные multilingual-эмбеддинги применимы для построения единого
        признакового пространства русско-английских почтовых сообщений
        и~могут использоваться как основа для кластеризации и~классификации.
  \item Кластеризация исторической почтовой выборки на~основе плотности
        позволяет формировать кандидаты на~персонализированные классы
        без~предварительного задания фиксированного числа категорий
        и~с~явным выделением outliers.
  \item Введение категории \texttt{other} и~порога уверенности снижает риск
        ошибочной автоматической классификации сообщений, не~принадлежащих
        известным классам, и~создаёт основу для управляемой обратной связи.
  \item Активация классификатора в~почтовом workflow должна выполняться
        только после прохождения quality gates, включающих метрики качества
        модели, анализ ошибок и~операционные показатели обработки сообщений.
\end{enumerate}

{\reliability} полученных результатов обеспечивается использованием открытых
методов машинного обучения и~обработки текста, воспроизводимой конфигурацией
экспериментов, фиксированными seed для повторных запусков, сохранением
артефактов моделей и~наборов данных, а~также сравнением с~базовыми подходами
к~классификации текстов. Программная реализация сопровождается автоматическими
проверками: unit-тестами, статическим анализом и~контролем типов.

\textbf{Структура работы.} Работа состоит из~введения, трёх глав, заключения,
списка использованных источников и~приложений.

Во~введении обосновывается актуальность темы, формулируются цель, задачи,
объект и~предмет исследования, методы, научная новизна, практическая
значимость, а~также положения, выносимые на~защиту.

В~первой главе анализируется предметная область автоматической организации
электронной почты, формулируется задача многоклассовой классификации
с~персонализированными классами и~категорией \texttt{other}, обзорно
рассматриваются методы представления текста, классификации и~кластеризации,
проводится сравнительный анализ существующих систем и~формализуются
функциональные и~нефункциональные требования к~разрабатываемой системе.

Во~второй главе описываются принципы проектирования и~общая архитектура
системы, модель данных, алгоритмы ingestion и~нормализации, алгоритм
формирования первичных классов на~основе кластеризации, архитектура
supervised-классификатора, стратегия принятия решений и~обработки категории
\texttt{other}, а~также критерии качества и~сценарий обратной связи через
переаннотацию сообщений.

В~третьей главе изложены технологический стек и~структура исходного кода,
реализация компонентов хранилища, IMAP-синхронизации, нормализации, очереди
и~воркеров, кластеризации, интеграции с~системой разметки и~обучения
классификатора, а~также приводятся результаты экспериментального
исследования: характеристика датасета, сравнение моделей эмбеддингов, оценка
кластеризации, обучение и~анализ классификатора, поведение системы
на~out-of-distribution входах и~операционные характеристики.

В~заключении формулируются основные результаты работы, подтверждается
достижение поставленных задач и~обозначаются направления дальнейшего развития.
Список использованных источников включает стандарты, монографии, статьи
и~документацию инструментов. В~приложениях приведены вспомогательные
материалы: диаграммы, конфигурации, листинги, таблицы и~протоколы
экспериментов.
